\documentclass{beamer}

\pdfmapfile{+sansmathaccent.map}


\mode<presentation>
{
	\usetheme{Warsaw} % or try Darmstadt, Madrid, Warsaw, Rochester, CambridgeUS, ...
	\usecolortheme{seahorse} % or try seahorse, beaver, crane, wolverine, ...
	\usefonttheme{serif}  % or try serif, structurebold, ...
	\setbeamertemplate{navigation symbols}{}
	\setbeamertemplate{caption}[numbered]
} 


%%%%%%%%%%%%%%%%%%%%%%%%%%%%
% itemize settings


%%%%%%%%%%%%%%%%%%%%%%%%%%%%
% itemize settings

\definecolor{myhotpink}{RGB}{255, 80, 200}
\definecolor{mywarmpink}{RGB}{255, 60, 160}
\definecolor{mylightpink}{RGB}{255, 80, 200}
\definecolor{mypink}{RGB}{255, 30, 80}
\definecolor{mydarkpink}{RGB}{155, 25, 60}

\definecolor{mypaleblue}{RGB}{240, 240, 255}
\definecolor{mylightblue}{RGB}{120, 150, 255}
\definecolor{myblue}{RGB}{90, 90, 255}
\definecolor{mygblue}{RGB}{70, 110, 240}
\definecolor{mycourseblue}{RGB}{120, 140, 255}
\definecolor{mycourselightblue}{RGB}{245, 250, 255}
\definecolor{mydarkblue}{RGB}{0, 0, 180}
\definecolor{myblackblue}{RGB}{40, 40, 120}

\definecolor{mygreen}{RGB}{0, 200, 0}
\definecolor{mygreen2}{RGB}{245, 255, 230}
\definecolor{mydarkgreen}{RGB}{0, 120, 0}
\definecolor{mydarkgreen}{RGB}{0, 120, 0}


\definecolor{mygray}{gray}{0.8}
\definecolor{mydarkgray}{RGB}{80, 80, 160}
\definecolor{mylightgray}{RGB}{160, 160, 160}

\definecolor{mydarkred}{RGB}{160, 30, 30}
\definecolor{mylightred}{RGB}{255, 150, 150}
\definecolor{myred}{RGB}{200, 110, 110}
\definecolor{myblackred}{RGB}{120, 40, 40}

\definecolor{mypink}{RGB}{255, 30, 80}
\definecolor{myhotpink}{RGB}{255, 80, 200}
\definecolor{mywarmpink}{RGB}{255, 60, 160}
\definecolor{mylightpink}{RGB}{255, 80, 200}
\definecolor{mydarkpink}{RGB}{155, 25, 60}
\definecolor{mywhitepink}{RGB}{255, 240, 240}

\definecolor{mydarkcolor}{RGB}{60, 25, 155}
\definecolor{mylightcolor}{RGB}{130, 180, 250}



\setbeamertemplate{itemize items}[default]

\setbeamertemplate{itemize item}{\color{myblackblue}$\blacksquare$}
\setbeamertemplate{itemize subitem}{\color{mygblue}$\blacktriangleright$}
\setbeamertemplate{itemize subsubitem}{\color{mygray}$\blacksquare$}

\setbeamercolor{palette quaternary}{fg=white,bg=mycourseblue}
\setbeamercolor{titlelike}{parent=palette quaternary}

\setbeamercolor{palette quaternary2}{fg=black,bg=mycourselightblue}
\setbeamercolor{frametitle}{parent=palette quaternary2}

\setbeamerfont{frametitle}{size=\Large,series=\scshape}
\setbeamerfont{framesubtitle}{size=\normalsize,series=\upshape}





%%%%%%%%%%%%%%%%%%%%%%%%%%%%
% block settings

\setbeamercolor{block title}{bg=red!30,fg=black}

\setbeamercolor*{block title example}{bg=mygreen!40!white,fg=black}

\setbeamercolor*{block body example}{fg= black, bg= mygreen2}


%%%%%%%%%%%%%%%%%%%%%%%%%%%%
% URL settings
\hypersetup{
	colorlinks=true,
	linkcolor=blue,
	filecolor=blue,      
	urlcolor=blue,
}

%%%%%%%%%%%%%%%%%%%%%%%%%%

\renewcommand{\familydefault}{\rmdefault}


\newcommand{\mycolor}[1] {
\setbeamercolor{normal text}{fg={#1}} % Recolors standard body text
\setbeamercolor{structure}{fg={#1}}   % Recolors bullets and structure items
\usebeamercolor[fg]{normal text}      % Applies the normal text color immediately
}


\usepackage{amsmath}
\usepackage{mathtools}

\usepackage{cancel}
\usepackage{centernot}

\usepackage{changepage}
\usepackage{subcaption}

\usepackage{qrcode}

\newcommand{\bo}[1] {\mathbf{#1}}
\newcommand{\R}{\mathbb{R}} 
\newcommand{\T}{^\top}     



\newcommand{\mydate}{Fall 2026}

\newcommand{\mygit}{\textcolor{blue}{\href{https://github.com/SergeiSa/Advanced-Control-2026}{github.com/SergeiSa/Advanced-Control-2026}}}

\newcommand{\myqr}{ \textcolor{black}{\qrcode[height=1.5in]{https://github.com/SergeiSa/Advanced-Control-2026}}
}

\newcommand{\myqrframe}{
	\begin{frame}
		\centerline{Lecture slides are available via Github:}
		\bigskip
		\centerline{\mygit}
		\bigskip
		\myqr
	\end{frame}
}


\newcommand{\bref}[2] {\textcolor{blue}{\href{#1}{#2}}}

%%%%%%%%%%%%%%%%%%%%%%%%%%%%
% code settings

\usepackage{listings}
\usepackage{color}
% \definecolor{mygreen}{rgb}{0,0.6,0}
% \definecolor{mygray}{rgb}{0.5,0.5,0.5}
\definecolor{mymauve}{rgb}{0.58,0,0.82}
\lstset{ 
	backgroundcolor=\color{white},   % choose the background color; you must add \usepackage{color} or \usepackage{xcolor}; should come as last argument
	basicstyle=\footnotesize,        % the size of the fonts that are used for the code
	breakatwhitespace=false,         % sets if automatic breaks should only happen at whitespace
	breaklines=true,                 % sets automatic line breaking
	captionpos=b,                    % sets the caption-position to bottom
	commentstyle=\color{mygreen},    % comment style
	deletekeywords={...},            % if you want to delete keywords from the given language
	escapeinside={\%*}{*)},          % if you want to add LaTeX within your code
	extendedchars=true,              % lets you use non-ASCII characters; for 8-bits encodings only, does not work with UTF-8
	firstnumber=0000,                % start line enumeration with line 0000
	frame=single,	                   % adds a frame around the code
	keepspaces=true,                 % keeps spaces in text, useful for keeping indentation of code (possibly needs columns=flexible)
	keywordstyle=\color{blue},       % keyword style
	language=Octave,                 % the language of the code
	morekeywords={*,...},            % if you want to add more keywords to the set
	numbers=left,                    % where to put the line-numbers; possible values are (none, left, right)
	numbersep=5pt,                   % how far the line-numbers are from the code
	numberstyle=\tiny\color{mygray}, % the style that is used for the line-numbers
	rulecolor=\color{black},         % if not set, the frame-color may be changed on line-breaks within not-black text (e.g. comments (green here))
	showspaces=false,                % show spaces everywhere adding particular underscores; it overrides 'showstringspaces'
	showstringspaces=false,          % underline spaces within strings only
	showtabs=false,                  % show tabs within strings adding particular underscores
	stepnumber=2,                    % the step between two line-numbers. If it's 1, each line will be numbered
	stringstyle=\color{mymauve},     % string literal style
	tabsize=2,	                   % sets default tabsize to 2 spaces
	title=\lstname                   % show the filename of files included with \lstinputlisting; also try caption instead of title
}


%%%%%%%%%%%%%%%%%%%%%%%%%%%%
% URL settings
\hypersetup{
	colorlinks=false,
	linkcolor=blue,
	filecolor=blue,      
	urlcolor=blue,
}

%%%%%%%%%%%%%%%%%%%%%%%%%%

%%%%%%%%%%%%%%%%%%%%%%%%%%%%
% tikz settings

\usepackage{tikz}
\tikzset{every picture/.style={line width=0.75pt}}
\usetikzlibrary{shapes.geometric, decorations.markings, positioning, patterns, calc, arrows.meta}


\subtitle{Advanced Control}
\author{by Sergei Savin}
\centering
\date{\mydate}



\title{Q Learning}
\begin{document}
\maketitle


\begin{frame}{Content}
	\begin{itemize}
		\item Model-based vs. model-free reinforcement learning
		\item Terminology: from optimal control to RL
		\item Return, rewards, and the discount factor
		\item Action-value function $Q^\pi(s,a)$ and $Q^*(s,a)$
		\item Bellman relation and the optimality equation
		\item Optimal policy from $Q^*$
		\item Temporal difference update (SARSA)
		\item Q-learning (off-policy update)
		\item Deep Q-learning (DQN)
		\item Literature
	\end{itemize}
\end{frame}


\begin{frame}{Reinforcement learning}
	\begin{flushleft}
		
		\textbf{Model-based optimal control}
		
		We know the system model:
		\[
		x_{k+1}=f(x_k,u_k,w_k).
		\]
		This allows us to compute (or approximate) value functions
		and derive an optimal control policy.
		
		\vspace{0.4cm}
		
		\textbf{Model-free reinforcement learning}
		
		We do not know the system model. Instead, we collect
		observations and learn directly from data.
		
		\vspace{0.3cm}
		
		\begin{itemize}
			\item We can learn value functions from observed transitions.
			\item Once we know the optimal value function, finding
			the optimal policy can be very simple.
		\end{itemize}
	\end{flushleft}
\end{frame}


\begin{frame}{Terminology: connection to optimal control}
	
	\centering
	\begin{tabular}{ll}
		\textbf{Optimal control} & \textbf{Reinforcement learning} \\
		\hline
		System & Environment \\
		Controller & Agent \\
		State $x$ & State $s$ \\
		Control input $u$ & Action $a$ \\
		Stage cost $l(x,u)$ & Reward $r(s,a)$ \\
		Cost $J$ (minimize) & Return $J$ (maximize) \\
		Control law & Policy $\pi$
	\end{tabular}
	
	\vspace{0.4cm}
	
	\begin{itemize}
		\item The agent interacts with the environment.
		\item The agent observes the state and chooses an action.
		\item The environment evolves and provides a reward.
	\end{itemize}
	
\end{frame}


\begin{frame}{Return function and rewards}
	\begin{flushleft}
	The reward function
	\[
	r_k = r(s_k,a_k)
	\]
	describes the immediate reward associated with the current
	state and action.
	
	\vspace{0.3cm}
	
	The discounted return from step $k$ is
	\[
	G_k =
	\sum_{i=k}^{\infty}
	\gamma^{i-k}r(s_i,a_i).
	\]
	
	where $0<\gamma<1$ is the discount factor.
	
	\vspace{0.3cm}
	
	\begin{itemize}
		\item Immediate rewards have the highest weight.
		\item Future rewards are progressively discounted.
		\item The goal is to maximize the expected return.
	\end{itemize}
	
	\end{flushleft}
\end{frame}


\begin{frame}{Action-value function}
	
	\begin{flushleft}
	The \textbf{on-policy action-value function} is
	\[
	Q^\pi(s,a)=
	\mathbb{E}_\pi\left[
	\sum_{k=0}^{\infty}\gamma^k r_{t+k}
	\;\middle|\;
	s_t=s,\ a_t=a
	\right].
	\]
	
	\vspace{0.2cm}
	
	It describes the expected discounted return when:
	\begin{itemize}
		\item We take action $a$ in state $s$ now.
		\item We follow policy $\pi$ afterwards.
	\end{itemize}
	
	\vspace{0.3cm}
	
	The \textbf{optimal action-value function} is
	\[
	Q^*(s,a)=\max_\pi Q^\pi(s,a).
	\]
	
	It describes the best expected discounted return achievable
	when taking action $a$ in state $s$ and following an
	optimal policy afterwards.
	
	\end{flushleft}
\end{frame}


\begin{frame}{Bellman relation}
	\begin{flushleft}
	
	We can decompose the discounted return:
	\[
	\begin{aligned}
		G_k
		&=r_k+\gamma r_{k+1}+\gamma^2r_{k+2}+\ldots\\
		&=r_k+\gamma G_{k+1}.
	\end{aligned}
	\]
	
	\vspace{0.3cm}
	
	Taking the expectation gives the Bellman equation.
	
	\textbf{Policy evaluation:}
	\[
	Q^\pi(s,a)=
	\mathbb{E}_\pi\left[
	r+\gamma Q^\pi(s',a')
	\right].
	\]
	
	\vspace{0.2cm}
	
	\textbf{Optimality equation:}
	\[
	Q^*(s,a)=
	\mathbb{E}\left[
	r+\gamma\max_{a'}Q^*(s',a')
	\right].
	\]
	
	\end{flushleft}
\end{frame}


\begin{frame}{Optimal policy}
	\begin{flushleft}
	
	If we know the optimal action-value function,
	we can obtain the optimal policy by selecting
	the action with the highest $Q$-value.
	
	\[
	\boxed{
		\pi^*(s)=\arg\max_a Q^*(s,a)
	}
	\]
	
	\vspace{0.4cm}
	
	In other words, we choose the action that promises the highest
			expected future return.

	
	\end{flushleft}
\end{frame}



\begin{frame}{Temporal Difference Update}
	\begin{flushleft}
	
	Recall the Bellman equation for policy evaluation:
	\[
	Q^\pi(s,a)=
	\mathbb{E}_\pi\left[
	r+\gamma Q^\pi(s',a')
	\right].
	\]
	
	\vspace{0.2cm}
	
	We do not know the expectation. Instead, we can use
	an observed transition $(s,a,r,s',a')$.
	
	\vspace{0.2cm}
	
	We update our estimate of the action-value function:
	\[
	\boxed{
		Q(s,a)\leftarrow Q(s,a)+\alpha
		\underbrace{
			\left[
			r+\gamma Q(s',a')-Q(s,a)
			\right]
		}_{\text{TD error}}
	}
	\]
	
	where $\alpha$ is the learning rate.
	
	\medskip
	
	This is called SARSA update (State → Action → Reward → State → Action). It is an \emph{on policy} method: $a'$ is determined by the policy.
	
	\end{flushleft}
\end{frame}


\begin{frame}{Q-learning}
	\begin{flushleft}
	
	Recall the optimality equation:
	\[
	Q^*(s,a)=
	\mathbb{E}\left[
	r+\gamma\max_{a'}Q^*(s',a')
	\right].
	\]
	
	\vspace{0.2cm}
	
	We can use an observed transition to update our estimate:
	\[
	\boxed{
		Q(s,a)\leftarrow Q(s,a)+\alpha
		\left[
		r+\gamma\max_{a'}Q(s',a')
		-Q(s,a)
		\right]
	}
	\]
	
	\vspace{0.3cm}
	
	This is the \textbf{Q-learning update}.  It is an \emph{off policy} method - we use the max over all actions to find $a'$, not policy.
	
	\vspace{0.2cm}
	
	\begin{itemize}
		\item We observe a transition $(s,a,r,s')$.
		\item We estimate the best future return using
		$\max_{a'}Q(s',a')$.
		\item We update our current estimate $Q(s,a)$.
	\end{itemize}
	
	\end{flushleft}
\end{frame}


\begin{frame}{Deep Q-learning}
	\begin{flushleft}
	
	In many problems, representing $Q(s,a)$ as a table is
	infeasible or undesirable.
	
	\vspace{0.3cm}
	
	Instead, we can approximate the action-value function
	using an artificial neural network (ANN):
	\[
	Q(s,a) \approx Q_\theta(s,a),
	\]
	where $\theta$ denotes the network parameters.
	
	\vspace{0.3cm}
	
	\begin{itemize}
		\item The ANN acts as a \textbf{critic}: it estimates
		the expected future return for each action.
		\item Instead of updating individual $Q$-values,
		we update the parameters $\theta$ of the network.
		\item This approach is called
		\textbf{Deep Q-learning} or
		\textbf{Deep Q-Network (DQN)}.
	\end{itemize}
	
	\end{flushleft}
\end{frame}


\begin{frame}{Literature}
	\begin{itemize}
		
		\item Mnih et al. (2013):
		\bref{https://arxiv.org/abs/1312.5602}
		{Playing Atari with Deep Reinforcement Learning. Original Deep Q-Network (DQN) paper.}
		
		\item Hugging Face:
		\bref{https://huggingface.co/learn/deep-rl-course/en/unit2/introduction}
		{Deep Reinforcement Learning Course, Unit 2: Introduction to Q-Learning.}
		
		\item Hugging Face:
		\bref{https://huggingface.co/learn/deep-rl-course/en/unit3/introduction}
		{Deep Reinforcement Learning Course, Unit 3: Deep Q-Learning and DQN.}
		
		\item PyTorch documentation:
		\bref{https://docs.pytorch.org/tutorials/intermediate/reinforcement_q_learning.html}
		{Reinforcement Learning (DQN) Tutorial: implementing DQN on the CartPole environment.}
		
		\item Stable-Baselines3:
		\bref{https://stable-baselines3.readthedocs.io/en/master/modules/dqn.html}
		{DQN documentation and reference implementation.}
		
	\end{itemize}
\end{frame}


\myqrframe






\end{document}
